\documentclass[a4paper,11pt]{article}
\usepackage[a4paper,margin=2cm]{geometry}
\usepackage[T1]{fontenc}
\usepackage[utf8]{inputenc}
\usepackage[ngerman,english]{babel}
\usepackage{lmodern,microtype,xcolor,parskip,enumitem}
\usepackage[most]{tcolorbox}
\usepackage{titlesec}
\usepackage[hidelinks]{hyperref}
\definecolor{HeaderBlueDark}{RGB}{70,110,170}
\definecolor{RowBlueLight}{RGB}{235,243,255}
\definecolor{RowBlueDark}{RGB}{220,233,250}
\definecolor{SoftGray}{RGB}{90,90,90}
\setlength{\parindent}{0pt}
\setlist[itemize]{leftmargin=1.4em,itemsep=0.35em,topsep=0.35em}
\titleformat{\section}[block]{\Large\bfseries\color{HeaderBlueDark}}{}{0pt}{}
\newtcolorbox{exbox}{enhanced,breakable,colback=RowBlueLight,colframe=RowBlueDark,boxrule=0.4pt,arc=1.5mm,left=1.2mm,right=1.2mm,top=1mm,bottom=1mm,title={Beispiele \textnormal{(Examples)}},fonttitle=\bfseries,colbacktitle=RowBlueDark,coltitle=black}
\NewDocumentEnvironment{grammarentry}{mm+b}{\begin{tcolorbox}[enhanced,breakable,colback=white,colframe=HeaderBlueDark,title={#1\hfill\normalfont\footnotesize #2},colbacktitle=HeaderBlueDark,coltitle=white,fonttitle=\bfseries]#3\end{tcolorbox}}{}
\newcommand{\GermanText}[1]{\textbf{Deutsch: }#1\par}
\newcommand{\EnglishText}[1]{{\small\color{SoftGray}\textbf{English: }#1\par}}
\newcommand{\ExampleItem}[2]{\item \textbf{#1}\\{\small\color{SoftGray}#2}}
\title{Die Modalpartikel \glqq mal\grqq\\{\large\textnormal{The modal particle mal}}}
\date{}
\begin{document}
\maketitle
\tableofcontents
\section{Bedeutung und Wirkung \textnormal{(Meaning and effect)}}
\begin{grammarentry}{mal}{Modalpartikel / Abtönungspartikel}
\GermanText{Mal kann Aufforderungen oder Bitten lockerer und weniger schroff wirken lassen.}
\EnglishText{Mal can make commands or requests sound more casual and less blunt.}
\end{grammarentry}
\section{Grammatische Rolle und Satzarten \textnormal{(Grammar and sentence moods)}}
\begin{grammarentry}{Verwendung}{Satzart und Kontext}
\GermanText{Besonders in Imperativen und höflichen Fragen. Mal garantiert nicht in jedem Kontext Höflichkeit.}
\EnglishText{Especially in imperatives and polite questions. Mal does not guarantee politeness in every context.}
\end{grammarentry}
\begin{grammarentry}{Form und Stellung}{Syntax}
\GermanText{Als Modalpartikel ist mal unflektierbar, fordert keinen Kasus und steht normalerweise im Mittelfeld. Allein im Vorfeld eines normalen Aussagesatzes kann eine Modalpartikel nicht stehen. Die Betonung und der Kontext beeinflussen die Wirkung.}
\EnglishText{As a modal particle, mal is invariable, does not govern case and normally appears in the middle field. It cannot occupy the prefield of a normal statement on its own. Stress and context affect its pragmatic force.}
\end{grammarentry}
\section{Typische Beispiele \textnormal{(Typical examples)}}
\begin{exbox}\begin{itemize}
\ExampleItem{Schau mal!}{Have a look!}
\ExampleItem{Komm mal her!}{Come over here for a moment!}
\ExampleItem{Kannst du mir mal helfen?}{Could you give me a hand?}
\ExampleItem{Ruf doch mal an!}{Why don't you call sometime!}
\end{itemize}\end{exbox}
\section{Abgrenzung zu anderen Verwendungen \textnormal{(Contrast with other uses)}}
\begin{grammarentry}{Wortarten unterscheiden}{Nicht immer Modalpartikel}
\GermanText{Mal als Temporaladverb heißt einmal oder irgendwann; in der Mathematik bedeutet mal multipliziert mit.}
\EnglishText{Temporal mal means once or sometime; in arithmetic mal means times.}
\end{grammarentry}
\begin{exbox}\begin{itemize}
\ExampleItem{Ich war mal in Wien.}{I was in Vienna once.}
\ExampleItem{Drei mal vier ist zwölf.}{Three times four is twelve.}
\end{itemize}\end{exbox}
\section{Lernhinweis \textnormal{(Learning tip)}}
\GermanText{Eine Modalpartikel verändert vor allem die Gesprächsabsicht. Versuchen Sie nicht, jedem Vorkommen eine feste englische Eins-zu-eins-Übersetzung zu geben. Achten Sie auf Satzart, Tonfall und Situation.}
\EnglishText{A modal particle mainly affects the speaker's intention. Do not try to give it one fixed English translation in every context. Pay attention to sentence mood, intonation and situation.}
\section{Quellen \textnormal{(Sources)}}
\begin{itemize}
\item IDS Mannheim: \url{https://grammis.ids-mannheim.de/terminologie/2}
\item Duden: \url{https://www.duden.de/sprachwissen/sprachratgeber/Modalpartikeln}
\end{itemize}
\end{document}
